\begin{tabularx}{\textwidth}{L{27mm}X X}
\toprule
\textbf{ER elem} & \textbf{Jelentés} & \textbf{Jelölés és példa} \\
\midrule
Egyed & Önálló léttel bíró objektumtípus, amelyről adatokat akarunk tárolni. & Téglalap; például \code{Hallgato}, \code{Targy}, \code{Rendeles}. \\
Gyenge egyed & Önálló kulccsal nem azonosítható egyed, amely egy másik egyedtől függ. & Kettős téglalap vagy jelölt kapcsolat; azonosításában a tulajdonos kulcsa is részt vesz. \\
Kapcsolat & Egyedpéldányok közötti asszociáció. & Rombusz; például hallgató \emph{felvesz} tárgyat. \\
Tulajdonság & Egyedhez vagy kapcsolathoz tartozó adat. & Ellipszis; például név, dátum, kredit. \\
Kulcstulajdonság & Az egyedpéldányt azonosító tulajdonság. & Aláhúzott attribútum; például \code{hallgatoId}. \\
Összetett tulajdonság & Több elemi részből álló tulajdonság. & Például cím = irányítószám, város, utca. \\
Többértékű tulajdonság & Egy egyedpéldányhoz több értéket is felvehet. & Kettős ellipszis; például több telefonszám. \\
Származtatott tulajdonság & Más adatokból számítható érték. & Szaggatott ellipszis; például életkor a születési dátumból. \\
\bottomrule
\end{tabularx}
